\clearpage
\fancyhead[L]{\small OBSERVATIONS OF MAGNETIC INCLINATION.}
\SourceLeaf{450}\pdfbookmark[0]{Observations of magnetic inclination in Göttingen.}{inclination-observations-title}
\begin{center}
{\Large\bfseries OBSERVATIONS\par\medskip OF MAGNETIC INCLINATION IN GÖTTINGEN.}
\par\bigskip
Resultate aus den Beobachtungen des magnetischen Vereins. 1841. II.
\end{center}
\begingroup
\renewcommand{\Article}[1]{\par\medskip\Needspace{4\baselineskip}\pdfbookmark[1]{#1.}{inclination-observations-art-#1}\centerline{\textbf{#1.}}\nobreak\smallskip}%
\Article{1}
The dip circle used for the observations reported here is by \textsc{Robinson}. It was the last instrument of this kind supplied by that distinguished instrument maker.

The vertical circle has a clear diameter of $241.169$ millimetres and is graduated at intervals of ten minutes. The distance between the inner ends of two graduation marks is therefore $0.351$ millimetres. The marks look very elegant even under considerable magnification in the microscope. From measurements on several of them I found their width to be $=0.024$ millimetres, so that one mark covers nearly $41$ seconds of arc.

The diameter of the horizontal circle, measured where the end of the index line meets the graduation marks, is $148$ millimetres. The divisions are in half degrees and the vernier gives single minutes. Only one reading is taken.

The degrees on the vertical circle are numbered from both ends of a horizontal diameter, upwards and downwards to $90$. This arrangement may perhaps seem convenient in ordinary observations, but it easily causes confusion when a deliberately weighted needle is used and consequently moves into another quadrant, or when observations are also made in a vertical plane perpendicular to the magnetic meridian. At any rate, in such ca%
\SourceLeaf{451}%
ses the arrangement makes it necessary to keep a somewhat more laborious and less clearly arranged record. I would therefore prefer a continuous graduation in the same direction from $0$ to $360^0$, or twice from $0$ to $180^0$. I have made it a habit, in the lower quadrant on the left or the upper quadrant on the right, to write down at once the supplement to $180^0$ in place of the engraved reading. All the readings in this paper are given in that form. On the horizontal circle the numbers run twice in the same direction from $0$ to $180^0$. It would be more natural and convenient to number continuously up to $360^0$, and I have set down the readings that occur here in that form.

On the spirit level a deflection of one millimetre corresponds to an inclination of $9$ seconds of arc.

\Article{2}
The instrument has four needles, which I distinguish by the numbers $1,2,3,4$. The last two have rotatable axes, an arrangement I shall return to below. Microscope measurements revealed no difference in thickness among all eight pivots. I found this thickness to be $=0.590$ millimetres. Needles $1$ and $2$ weigh $16.5$ grams each, and the other two $20.5$ grams each.

There are small differences in the lengths of the individual needles. The measurements give
\begin{center}
\begin{tabular}{crcrl}
for &$1$&$.\quad.\quad.\quad.\quad.\quad.$&$240.931$&millimetres\\
&$2$&$.\quad.\quad.\quad.\quad.\quad.$&$240.866$&—\\
&$3$&$.\quad.\quad.\quad.\quad.\quad.$&$240.938$&—\\
&$4$&$.\quad.\quad.\quad.\quad.\quad.$&$240.954$&—
\end{tabular}
\end{center}
Thus the shortest needle is only $0.303$ millimetres shorter than the clear diameter of the circle, and the longest only $0.215$ millimetres shorter. This helps make the readings more precise, but it also means that even a very small eccentricity can impede the free motion of the needle. It is therefore difficult to adjust the parts of the instrument whose positions determine the eccentricity in an entirely satisfactory way, especially since their positions must also meet four other conditions, making \textit{six} in all.

\SourceLeaf{452}\Article{3}
These six conditions are as follows.

The two agate plates, on whose upper edges the pivots of the needle rest during observations, must be adjusted by the two pairs of screws that support them so that

1) their upper edges lie in a single plane,

2) this plane is perpendicular to the plane of the vertical circle, and

3) it lies below the centre of that circle at a distance equal to half the thickness of a pivot,

4) the line of intersection of those two planes makes a right angle with the vertical axis.

The cups by which the needle is lifted off the agate plates and set back on them, and which are screwed to the lifting frame with some freedom to shift. must also be adjusted so that after the needle is set down

5) its axis is perpendicular to the line of intersection last mentioned (in 4), and hence, in conjunction with condition 2, also perpendicular to the plane of the vertical circle, and at the same time

6) it meets the vertical diameter of the circle.

Conditions 1, 2, 4 together take the place of a single condition. When the upright axis of rotation is exactly vertical, a horizontal plane should touch the edges of both agate plates along their lengths, or along two lines. This assumes that the plane of the vertical circle is parallel to that axis of rotation and thus is also vertical with it. This can be regarded as the seventh condition. It is usually assumed tacitly, in reliance on the skill of the instrument maker. The instrument as it stands provides no means of testing it or, if necessary, correcting it.

\Article{4}
In making the observations to be given in this paper, I proceeded as follows to test the stated conditions, for lack of other means.

\SourceLeaf{453}To test the \textit{first} condition I used the flat glass of a so-called artificial horizon. After the frame with the cups had been removed, the glass was placed on the agate bearings with its ground, matt rear face turned upwards. If the condition is unmet, only \textit{one} agate plate will be touched along its whole length and the other at one endpoint. One can see this. if the error is not very small, with the naked eye. A spirit level placed on the glass gives greater precision and certainty by showing whether the glass has two different positions of contact or just one. It is easy to see that once the first condition is met, this spirit level can also be used to test the \textit{second} and \textit{fourth}.

To test the \textit{fifth} condition, the needle must be brought into two different positions of equilibrium. In both, the pivots have the same position on the bearings (or the same face of the needle is at the front), and the needle has only a moderate inclination to the horizontal line, but the end that was on the left in one position is on the right in the other. The most convenient way to obtain these two positions is to weight the needle suitably. The condition in question is met if the edge at each end of the needle projects just as far in front of the rim of the circle in one position as in the other. In regions where the inclination is only moderate, the two positions could already be obtained simply by turning the instrument through half a revolution, so that the plane of the circle is near the magnetic meridian in both cases.

A similar test can also be used for the \textit{second} condition, except that two opposite, nearly vertical positions of the needle must be produced. One arises on its own when the plane of the circle is brought nearly perpendicular to the magnetic meridian, and the other is obtained either by suitable weighting or by reversing the poles. It is easy to see, however, that this procedure is equivalent to the one mentioned above using the spirit level only if the seventh condition is met. Combining both methods thus gives a kind of test of that condition itself. It can, of course, be only very imperfect, since one can judge only by estimation whether the edge of the needle projects equally far in front of the rim of the circle.

\expandafter\def\csname HistoricalNoteText454\endcsname{Except for those of 23 September.}%
\SourceLeaf{454}The same combined positions of the needle also serve to test the other two conditions. The \textit{sixth} condition is met if each end of the needle is as far from the inner face of the circle in the first nearly horizontal position as in the second. The same applies to the \textit{third} condition in the nearly vertical positions. If the two halves of the needle were exactly equal in length, it would clearly suffice to compare the distances of both ends of the needle from the inner face of the circle with one another in one nearly horizontal and one nearly vertical position. But with our instrument, where the clearances are so very small in general, this is insufficient, and even a very slight inequality between the two halves of the needle becomes apparent.

\Article{5}
How difficult it is to meet all the conditions at once in this way is clear from the fact that the two screws on which each agate plate rests are only eight millimetres apart. Since the pitch of a screw thread is $0.283$ millimetres, half a turn of one screw already turns the corresponding agate plate by one degree.

The task is made much easier by a special device that I had made only later, which brings the edges of the agate plates into a single plane and makes that plane horizontal. I shall pass over its description for now, since it could not yet be used for the observations reported here\HistoricalNote{454}. A second device, likewise completed only after the observations ended, serves to determine precisely the deviation of the main circle from the vertical. It gave a deviation of \textit{ten minutes}, but removing the deviation will first require an alteration to the instrument. In any event, its effect on the inclinations cannot amount even to one second of arc.

In general, I must point out that small errors in the various adjustments can have only a barely perceptible effect on the determination of inclina%
\SourceLeaf{455}%
tion. The effect that some of the errors have on the position of the needle is only of second order in those errors. The effect of the others, notably eccentricity and a tilt of the plane touching the agate plates in the direction parallel to the plane of the circle (errors with respect to conditions 3, 6 and 4), is completely eliminated by combining the individual observations. I therefore cannot agree with \textsc{Horner}'s view that one should attend \textit{above all} to removing this last error (Physik. Wörterb., vol. 5, p. 759). I regard it as the error whose complete removal matters least.

\Article{6}
The inclination observations to be given here were all made outdoors at the site described in the \textit{Resultate} 1840. II. [p. 433 of this volume]. A shelter kept the sun's rays off the instrument. It was set up on the stone so that the straight line through the points of two feet was nearly perpendicular to the magnetic meridian. The places of all three feet were marked for that position. The instrument was oriented precisely with respect to the magnetic meridian by an auxiliary needle supplied with it, which is suspended from a point by means of a small agate cap. The support for that point has two short cylindrical side arms that are placed in the two cups. The point then stands vertical of its own accord under the weight of the freely hanging part of the support. I have often used this method of orientation together with the usual method of taking corresponding inclinations in two positions of the vertical circle nearly perpendicular to the magnetic meridian. I have always found only quite negligible differences, which shows that the auxiliary needle is sufficiently sensitive and introduces no constant deviation. Incidentally, a slight departure of the vertical plane in which one observes from the magnetic meridian, which is itself not entirely constant during the observations, has only a second-order effect on the tilt of the dip needle, which can be regarded as quite imperceptible.

\SourceLeaf{456}\Article{7}
Even the most skilled instrument makers can make the centre of gravity of a needle coincide with the axis of rotation only approximately. There is almost always some remaining offset. Its effect on the position in which the needle settles is to be determined or eliminated by combining observations made under several different sets of conditions. Reversing the poles of the needle is an essential part of varying those conditions. Other things being equal, the more weakly the needle is magnetised, the stronger that effect. But one is not entitled to assume. that the needle will have the same magnetic strength after its poles have been reversed as it did before. An exact reduction of the observations therefore depends on knowing the ratio of these quantities. This can be found by observing the swing duration of the needle. For several reasons, however, I preferred \textit{horizontal} swings, and to observe them I used a special apparatus made by Inspector \textsc{Meyerstein}. The needle swings in a wooden box with glazed lids. It rests on a light yoke suspended by a silk thread $270$ millimetres long, protected from draughts by a glass tube. As the needle swings, its ends move along two graduated arcs, each spanning $40$ degrees and divided into half degrees, so that $5$ minutes can be estimated with confidence. Each needle's swing duration was determined before and after it was stroked to reverse its poles, each time from $150$ swings divided into three sets. After proper reduction to infinitely small arcs, the sets always give results in excellent agreement. The observations usually began with an arc of swing of about $36$ degrees. It is worth noting here that, in contrast to the experience with heavier bars mentioned in the \textit{Resultate} 1837. IV [p. 385 of this volume], the arc of swing decreased at nearly the same rate on all days and for all needles. The time in which the arc fell to a quarter of its original value was thus $14$ minutes, with slight variations. These observations of swings were always made on a stone pedestal in the observatory, since the ratio matters here more than the absolute duration. The small external influences that might be present there cannot affect the ratio appreciably.

\expandafter\def\csname HistoricalNoteText457\endcsname{The most striking example of this kind is given on p. 142 of the \textit{Fifth Report of the British association for the advancement of Science}. The eight needles with which Captain Ross determined the inclination in London gave differences of up to 41 minutes, even though the observations with each individual needle were numerous and agreed well with one another. In England this unusual occurrence, about which no further details are given, was attributed to the pivots not being perfectly cylindrical, and it was for this very reason that rotatable pivots were tried.}%
\SourceLeaf{457}\Article{8}
In the inclination observations made in the summer of 1842, I in\-tended both to establish the magnetic inclination for that time and to determine what degree of accuracy can be achieved with the instrument in use. It seemed insufficient to me to estimate the reliability of the final results, which are influenced by many different circumstances, from the differences in the needle's settling positions when it is repeatedly lifted \textit{by means of the cup lever}. Nor can that purpose be served by merely comparing the inclination results obtained from observations on different days, since the random observational errors attributable to the instrument are then mixed with actual fluctuations in the inclination itself. I was also eager to find out whether my four needles would give results that agreed, or whether they would give distinctly and substantially different results, as has happened to some observers\HistoricalNote{457}.

These considerations led me to arrange the observations somewhat differently from usual. The essential difference will be clear from what follows.

The position of the needle, that is, the position of both tips with respect to the graduations on the circle, is usually observed in four different combinations of the position of the circle and the way the needle is placed in it. The graduated face of the circle and the marked face of the needle can face east or west, in the same or in opposite directions. The same combinations are repeated after reversing the poles, giving $16$ readings in all. Their simple arithmetic mean is taken as the inclination provided they do not differ greatly as a result of a large offset of the needle's centre of gravity from its pivot axis. In the opposite case a more elaborate calculation is used. It goes without saying that each of the $16$ readings
\SourceLeaf{458}may itself already be the mean of a smaller or larger number of settling positions obtained by repeatedly lifting the needle in each combination.

My procedure differs in that I observed with two needles on each day, without reversing their poles between the observations. The poles were reversed between two successive observations, and always on just one needle, alternating between the two. It is clear that in this way the observations of four days cover all combinations of the different polarisations of the two needles, as happened with needles 1 and 3 from 6 to 9 July and with needles 2 and 4 from 17 to 20 July. Continuing these similarly combined alternations over eight days of observation, as was done with needles 1 and 2 from 20 May to 5 June and with needles 3 and 4 from 8 to 25 June, thus gave each combination of polarisations twice. Each day's observations were arranged so that the results from both needles would be as nearly simultaneous as possible. To achieve this, the four combinations mentioned above were first observed in full with one needle, each with the needle laid back four times. The corresponding combinations were then observed with the second needle, with the needle laid back eight times. Finally the same combinations were observed again with the first needle, but in reverse order and with the needle laid back four times.

With this arrangement, a single day's observations by themselves do not yet give a determination of the inclination. But if they are combined with the following day's observations, the needle whose magnetisation was left unchanged clearly shows how much the inclination differed on the two days. The observations made with the other needle in only one polarisation can then be reduced to a single time and thus made complete. For a more rigorous treatment covering all the observations of all 24 days, however, the relation of the partial results to one another must first be discussed in more detail.

\Article{9}
The most convenient way to begin this development, which is important in more than one respect, is with an example taken from an observation made in the cus%
\SourceLeaf{459}%
tomary way. I made such observations on several days as well.

For this purpose I choose the observation with needle 1 on the morning of 23 September 1842, from $8\frac{1}{2}$ to $11$ o'clock. The magnetic orientation was obtained with the auxiliary needle in the way described in Article 6. The index of the azimuth circle, whose degrees increase from left to right and which I count from $0$ to $360^0$, as noted above, indicated $90^0\ 5'$ when the vertical circle was in the meridian with its graduated side facing east.

Besides the usual eight combinations in the magnetic meridian, I made just as many in the vertical plane perpendicular to it on that day. I include these observations here, since they give occasion for several matters to discuss. The needle, like the other three, is marked on one side with the letters $A$, $B$ at its ends, so that the polarisation and the way it is placed can easily be distinguished. In each of the $16$ combinations the needle was laid on the agate plates five times by means of the cup lever. In the following overview, however, I give only the mean values of the corresponding settling positions.

\begin{center}
\textit{Needle end $B$ north pole.}\par\smallskip
{\fontsize{9}{11}\selectfont
\begin{tabular}{c||cc||cc}
Azimuth&\multicolumn{4}{c}{Marked face of the needle}\\
circle&\multicolumn{2}{c||}{front}&\multicolumn{2}{c}{back}\\
&below&above&below&above\\\hline
$90^0\ 5'$&$67^0\ 27'\ 54''$&$67^0\ 29'\ 36''$&$67^0\ 45'\ 39''$&$67^0\ 44'\ 51''$\\
$180\ 5$&$89\ 52\ 39$&$89\ 52\ 51$&$90\ 12\ 30$&$90\ 10\ 30$\\
$270\ 5$&$112\ 18\ 39$&$112\ 16\ 45$&$112\ 38\ 51$&$112.\ 33\ 54$\\
$0\ 5$&$89\ 58\ 33$&$89\ 57\ 48$&$90\ 13\ 27$&$90\ 10\ 54$
\end{tabular}}
\par\medskip
\textit{Needle end $A$ north pole.}\par\smallskip
{\fontsize{9}{11}\selectfont
\begin{tabular}{c||cc||cc}
$90^0\ 5'$&$68^0\ 2'\ 51''$&$68^0\ 2'\ 33''$&$67^0\ 35'\ 15''$&$67^0\ 37'\ 0''$\\
$180\ 5$&$90\ 14\ 48$&$90\ 12\ 21$&$89\ 51\ 12$&$89\ 51\ 36$\\
$270\ 5$&$112\ 27\ 21$&$112\ 22\ 33$&$112\ 7\ 6$&$112\ 5\ 33$\\
$0\ 5$&$90\ 16\ 15$&$90\ 14\ 0$&$89\ 53\ 54$&$89\ 54\ 18$
\end{tabular}}
\end{center}

\Needspace{6\baselineskip}
\SourceLeaf{460}The duration of a horizontal swing was found to be
\begin{center}
\begin{tabular}{lll}
before the observations&$.\quad.\quad.\quad.\quad.\quad.\quad.\quad.\quad.$&$5''83555$\\
after the observations&$.\quad.\quad.\quad.\quad.\quad.\quad.\quad.\quad.$&$5.87416$
\end{tabular}
\end{center}

\Article{10}
I turn first to the differences between the readings of the lower and upper tips. These differences depend on the fact that the pivot axis passes neither through the centre of the graduations nor through the straight line joining the two tips of the needle. Let $x$, $y$ denote the coordinates of the intersection of the pivot axis with the plane of the circle, measured from the centre of the graduations in divisions of arc along the inner circumference. Let $x$ be parallel to the diameter through the two zero points and positive to the right, and $y$ parallel to the diameter through the two $90^0$ points and positive upwards. Further, let $180^0-z$ be the angle between the two planes drawn through the pivot axis and the tips $A$ and $B$, with the convention that, if the needle is imagined horizontal with its marked side facing upwards, the angle is counted from $A$ to $B$ from left to right. Finally, let $l$ be the mean of the two readings. Their difference, understood as the lower reading subtracted from the upper, is then
\[
=2x\sin l+2y\cos l\pm z
\]
where the upper sign applies when the marked face is at the front and the upper tip is $A$ (the south pole here), or when that face is at the back and $B$ is uppermost. The lower sign applies in the other two cases.

The observations above thus give $16$ equations, from which the method of least squares gives
\[
\begin{aligned}
x&=-38''3\\
y&=+153.2\\
z&=+75.4
\end{aligned}
\]
\Needspace{20\baselineskip}
The comparison then gives the following when the results are arranged by the magnitude of $l$
\SourceLeaf{461}\begin{center}
{\fontsize{9}{11}\selectfont
\begin{tabular}{c|r|r|r}
$l$&Observation&Calculation&Error\\\hline
$67^0\ 28'\ 45''$&$+102''$&$+122''$&$-20''$\\
$67\ 36\ 7$&$+105$&$+121$&$-16$\\
$67\ 45\ 15$&$-48$&$-30$&$-18$\\
$68\ 2\ 42$&$-18$&$-32$&$+14$\\
$89\ 51\ 24$&$+24$&$0$&$+24$\\
$89\ 52\ 45$&$+12$&$-1$&$+13$\\
$89\ 54\ 6$&$+24$&$-1$&$+25$\\
$89\ 58\ 10$&$-45$&$-1$&$-44$\\
$90\ 11\ 30$&$-120$&$-153$&$+33$\\
$90\ 12\ 10$&$-153$&$-153$&$0$\\
$90\ 13\ 34$&$-147$&$-153$&$+6$\\
$90\ 15\ 7$&$-135$&$-153$&$+18$\\
$112\ 6\ 19$&$-93$&$-111$&$+18$\\
$112\ 17\ 42$&$-114$&$-112$&$-2$\\
$112\ 24\ 57$&$-288$&$-264$&$-24$\\
$112\ 36\ 22$&$-297$&$-264$&$-33$
\end{tabular}}
\end{center}

The sum of the squares of the remaining errors is $7924$, from which it is inferred that the mean error of the difference of two means of five readings can be taken as
\[
=\BareRoot{\frac{7924}{13}}=24''7
\]
and the mean error of a single reading of one tip as
\[
=\BareRoot{\frac{5}{2}.\frac{7924}{13}}=39''0
\]
This is indeed a very satisfactory accuracy, which a similar analysis of the observations from other days confirms and sometimes even surpasses. It may be worth noting, however, that achieving such agreement depends essentially on lifting the needle only when it is at rest or close to its resting position. Without this precaution, the needle, whose swing consists in the pivots rolling on the bearing, would come to rest at a different place on the bearing from where it is laid down. The eccentricity component $x$ would then be variable.

The method described here does yield the values of the eccentricity components $x$ and $y$ with considerable accuracy, but these values
\SourceLeaf{462}cannot directly tell us whether, and by how much, the bearings and cups still have to be moved to meet conditions 3 and 6 in Article 3. The conditions refer to the centre of the inner circle, whereas the values refer to the centre of the graduations. These centres may differ slightly, and they actually do differ in the instrument in question. In fact, before the observations given here, the relevant adjustments had been made with every possible care.

The quantity denoted by $z$ is evidently constant for each needle, and a similar treatment of the observations from other days has given nearly the same value For the other three needles I found
\begin{center}
\begin{tabular}{crll}
for needle&$2$&$.\quad.\quad.\quad.\quad.\quad.\quad.\quad.$&$+3'\ 18''$\\
&$3$&$.\quad.\quad.\quad.\quad.\quad.\quad.\quad.$&$-1\ 4$\\
&$4$&$.\quad.\quad.\quad.\quad.\quad.\quad.\quad.$&$+1\ 2$
\end{tabular}
\end{center}

Although knowing these values is of no particular practical interest, their smallness does credit to the care the distinguished instrument maker devoted to making the needles.

\Article{11}
The mean of the readings at the two tips gives us the inclination of the straight line joining those tips, or of a line parallel to it, to the diameter of the vertical circle marked $0$. I set out these $16$ means here in pairs.
\begin{center}
\textit{Needle end $B$ north pole}\par\smallskip
{\fontsize{9}{11}\selectfont
\begin{tabular}{c|c||c|c}
Azim.&Marked needle face&Azim.&Marked needle face\\
circle&front&circle&back\\\hline
$90^0\ 5'$&$67^0\ 28'\ 45''$&$270^0\ 5'$&$112^0\ 36'\ 23''$\\
$180\ 5$&$89\ 52\ 45$&$0\ 5$&$90\ 12\ 11$\\
$270\ 5$&$112\ 17\ 43$&$90\ 5$&$67\ 45\ 15$\\
$0\ 5$&$89\ 58\ 10$&$180\ 5$&$90\ 11\ 30$
\end{tabular}}
\par\medskip
\textit{Needle end $A$ north pole}\par\smallskip
{\fontsize{9}{11}\selectfont
\begin{tabular}{c|c||c|c}
$90^0\ 5'$&$68^0\ 2'\ 42''$&$270^0\ 5'$&$112^0\ 6'\ 20''$\\
$180\ 5$&$90\ 13\ 34$&$0\ 5$&$89\ 54\ 6$\\
$270\ 5$&$112\ 24\ 57$&$90\ 5$&$67\ 36\ 7$\\
$0\ 5$&$90\ 15\ 8$&$180\ 5$&$89\ 51\ 24$
\end{tabular}}
\end{center}

\SourceLeaf{463}The settings placed side by side here are those in which the pivot axis had the same orientation relative to the points of the compass while the vertical circle was in opposite positions. The relation between two such numbers $l$ and $l'$ is very simple if the bearings have been adjusted so that a plane perpendicular to the vertical axis of rotation touches them. On that assumption the pivot axis lies in a horizontal plane in both settings, so the needle plainly comes to rest in the same position. That is, if by $L$ we mean the inclination of the straight line drawn from the upper tip to the lower tip to the horizontal radius of the circle that lies on the right of the marked face of the needle in each setting, then $L$ has the same value in both settings. The inclination is given by
\[
\begin{aligned}
\text{in the first setting }&=l-\alpha\\
\text{and in the second }&=180^0-(l'-\alpha)
\end{aligned}
\]
if $\alpha$ denotes the error of the zero point, that is, the reading at the radius of the circle that makes a right angle with the vertical axis. Under the assumption above we therefore have
\[
\alpha=\tfrac{1}{2}(l+l')-90^0,\qquad L=\tfrac{1}{2}(l+180^0-l')
\]
From the observations on 23 September, when this adjustment had been made with the greatest care using the device mentioned in Article 5, we thus obtain eight different determinations of $\alpha$
\begin{center}
\begin{tabular}{r}
$+2'\ 34''$\\
$2\ 28$\\
$1\ 29$\\
$4\ 50$\\
$4\ 31$\\
$3\ 50$\\
$0\ 32$\\
$3\ 16$
\end{tabular}
\end{center}

The sum of the squares of the deviations, expressed in seconds, from the mean $2'56''$ is found to be $=57214$. If these deviations are treated as entirely random, they therefore give the mean error of the
\SourceLeaf{464}result from a pair of corresponding settings as $=\BareRoot{\frac{57214}{7}}=90''4$. We see that with this instrument the anomalies in the settings are much larger than the errors of reading alone.

\Article{12}
The situation is different, however, if the bearings have not been adjusted as precisely as assumed. Suppose that their edges do lie in a single plane, but that this plane is not perpendicular to the vertical axis. The plane then tilts in opposite directions relative to the horizontal plane in the two settings. Only its tilt along the bearing edges, or parallel to the plane of the circle, matters here, since a slight tilt across them, or along the axis of the needle, has no perceptible effect on the position in which the needle comes to rest. Let $L$ denote the inclination of the needle, in the same sense as above, that would occur if it rested on a perfectly horizontal bearing. Let $\delta$ denote the corresponding directive force, that is, the coefficient by which the sine of a deflection from the resting position must be multiplied to express the torque of the force driving the needle back to that position. Finally, let $L+\beta$ be the inclination that actually occurs on the tilted bearing in the first setting. It is then easy to show that
\[
\delta\sin\beta=p\rho\sin\gamma
\]
where $p$ denotes the weight of the needle, $\rho$ the radius of the pivots, and $\gamma$ the tilt of the bearing to the horizontal. The last quantity is taken positive when the bearing is lower on the right of the marked face of the needle. In the second setting, however, we must evidently put $-\gamma$ in place of $\gamma$, which changes $\beta$ to $-\beta$, so the inclination of the needle in this second setting becomes $L-\beta$. We therefore have
\[
l-\alpha=L+\beta,\qquad180^0-(l'-\alpha)=L-\beta
\]
and hence, just as in the preceding article,
\[
\tfrac{1}{2}(l+180^0-l')=L
\]
while in place of the other equation there we now have
\[
\tfrac{1}{2}(l+l')-90^0=\alpha+\beta
\]


\expandafter\def\csname HistoricalNoteText465\endcsname{Strictly speaking, it is the sum of the moments of inertia of the needle and the yoke. Separating the two is partly impracticable and partly unnecessary, since this yoke is used only for horizontal swings.}
\SourceLeaf{465}But if the edges of the agate plates do not lie in a single plane at all, these same two formulas will still be sufficiently accurate if we take $\gamma$ to be the mean of the inclinations of the two edges, provided that the centre of gravity of the needle is nearly equidistant from the two bearing points of the pivots. Strictly speaking, a small modification is still required because the straight line joining the two points of contact between the pivots and the bearings does not have quite the same azimuth in the two settings. But the effect of this circumstance on the position of the needle may be regarded as quite imperceptible even where it is strongest, namely in observations in the plane perpendicular to the magnetic meridian.

\Article{13}
Since it is of some interest to see how much the tilt of the bearings affects the setting of the needle when the bearings have not been adjusted, I add here what is needed to determine this for the needle used on 23 September. To determine its moment of inertia, I had previously observed its horizontal swings both without and with a ring resting on it. The moment of inertia of the ring itself could be calculated with sufficient precision from its weight and dimensions. On 21 September the figures were
\begin{center}
{\fontsize{9}{11}\selectfont
\begin{tabular}{llrl}
Swing duration without ring&$.\quad.\quad.\quad.$&$5''88431$&\\
\hspace{2em}---\quad with ring&$.\quad.\quad.\quad.$&$7.32835$&\\
Weight of the ring&$.\quad.\quad.\quad.\quad.\quad.\quad.$&$19.2385$&grams\\
Inner diameter&$.\quad.\quad.\quad.\quad.\quad.$&$75.525$&millimetres\\
Outer diameter&$.\quad.\quad.\quad.\quad.\quad.$&$79.767$&---
\end{tabular}}
\end{center}

\Needspace{6\baselineskip}
Taking the gram and millimetre as units, we obtain
\begin{center}
\begin{tabular}{llr}
Moment of inertia of the ring&$.\quad.\quad.\quad.$&$29019$\\
\hspace{2em}---\quad of the needle\footnotemark&$.\quad.\quad.\quad.$&$52662$
\end{tabular}
\end{center}
\footnotetext{\csname HistoricalNoteText465\endcsname}


\SourceLeaf{466}Combining this with the swing durations given above in Article 9 for 23 September, and taking the length of the simple seconds pendulum at Göttingen to be $994.126$ millimetres, we obtain in the usual way
\begin{center}
\begin{tabular}{llr}
\multicolumn{3}{l}{horizontal magnetic directive force}\\
\quad before restroking&$.\quad.\quad.\quad.$&$1.5556$\\
\quad after restroking&$.\quad.\quad.\quad.$&$1.5352$
\end{tabular}
\end{center}

Strictly speaking. these numbers apply in the first instance only to the place where the swings were observed and thus include any local influences present there. This effect, which in any case is only small, is immaterial for the present purpose.

With an inclination of $67^0\ 40'\ 54''$ we further obtain
\begin{center}
\begin{tabular}{llr}
\multicolumn{3}{l}{total magnetic directive force,}\\
\quad before restroking&$.\quad.\quad.\quad.$&$4.0965$\\
\quad after restroking&$.\quad.\quad.\quad.$&$4.0429$\\
\multicolumn{3}{l}{vertical magnetic directive force}\\
\quad before restroking&$.\quad.\quad.\quad.$&$3.7897$\\
\quad after restroking&$.\quad.\quad.\quad.$&$3.7401$
\end{tabular}
\end{center}

If we disregard the small modification of the needle's magnetic directive force caused by the eccentricity of its centre of gravity, these four numbers can be taken as the values of $\delta$, according as the observation was made in the magnetic meridian or in the plane perpendicular to it. Since $\beta$ and $\gamma$ are always small enough that these quantities themselves can be used in place of their sines, so that
\[
\beta=\frac{p\rho}{\delta}.\gamma
\]
we obtain the following, according as the strength of magnetisation before or after restroking is taken as the basis.

For observations in the magnetic meridian
\[
\beta=1{,}1882\gamma\quad\text{or}\quad\beta=1{,}2039\gamma
\]
For observations in the plane perpendicular to the magnetic meridian
\[
\beta=1{,}2844\gamma\quad\text{or}\quad\beta=1{,}3014\gamma
\]


\SourceLeaf{467}The relations between the individual observations considered so far are of course not essential if all we want is to derive the magnetic inclination from them all. But they are of some importance in testing and confirming the result, since genuine confidence in the whole can arise only from a clear understanding of the satisfactory agreement among the parts.

\Article{14}
The observations have now been reduced to the eight values of $L$. These can be understood as the inclinations, in equilibrium, of the straight line drawn from the south pole tip of the needle to the north pole tip, to the horizontal radius of the circle on the right of the marked face of the needle. This assumes that the pivots rest on a horizontal surface, or, what is evidently the same from the standpoint of statics, that the needle can turn only about the axis of its pivots. In other words, the values of $L$ are the corrected values of the numbers listed in Article 11 under the heading \textit{Marked needle face at the front}. They are free of the effects of the zero-point error and the departure of the bearings from the horizontal
\begin{center}
\textit{Values of $L$.}\par\smallskip
\begin{tabular}{c|c|c}
Az. circle&$B$ north pole&$A$ north pole\\\hline
$90^0\ 5'$&$67^0\ 26'\ 11''$&$67^0\ 58'\ 11''$\\
$180\ 5$&$89\ 50\ 17$&$90\ 9\ 44$\\
$270\ 5$&$112\ 16\ 14$&$112\ 24\ 25$\\
$0\ 5$&$89\ 53\ 20$&$90\ 11\ 52$
\end{tabular}
\end{center}

To express in an equation how the values of $L$ depend on the elements that determine them, I use the following notation.

$V$\quad Position of the azimuth circle for the observation.

$V^0$\quad Position of the azimuth circle at which the vertical circle is in the magnetic meridian and its graduated side faces east.

$i$\quad Magnetic inclination.


\SourceLeaf{468}$m$\quad The product of the magnetic moment of the needle and the total intensity of the Earth's magnetic force, taking gravity as the unit of accelerating forces.

$q$\quad The weight of the needle multiplied by the distance of its centre of gravity from the pivot axis.

$c$\quad The acute angle between the straight line joining the tips of the needle and its magnetic axis, positive if the latter lies to the right when the needle is imagined horizontal with its marked face upwards.

$Q$\quad The angle between the straight line from the south pole tip of the needle to the north pole tip and the straight line from the pivot axis to the centre of gravity, reckoned from the first line, from left to right, with the needle in the same position as for $c$.

$\delta$\quad The directive force.

Resolving the Earth's magnetic force into a vertical and a horizontal component, we get from the first the torque, taken positive in the direction of increasing $L$,
\[
m\sin i\cos(L+c)
\]
and from the other
\[
-m\cos i\cos(V-V^0)\sin(L+c)
\]
Gravity, on the other hand, produces the torque
\[
q\cos(L+Q)
\]
Since $L$ specifies the equilibrium position, the sum of these three torques is $=0$. This gives us the main equation
\[
-\sin i\cos(L+c)+\cos i\cos(V-V^0)\sin(L+c)=\frac{q}{m}.\cos(L+Q)
\]
If we write $L+z$ in place of $L$ in the sum of the three torques, we get the torque that occurs at a deflection $z$ from equilibrium. If this expression is expanded in two parts with the factors $\cos z$ and $\sin z$, the first vanishes by the main equation, and the second is $=-\delta\sin z$ by the definition of the directive force.
\SourceLeaf{469}We thus have the general formula for $\delta$
\[
\delta=m\sin i\sin(L+c)+m\cos i\cos(V-V^0)\cos(L+c)+q\sin(L+Q)
\]
For the three principal cases this gives the following.

I.\quad For $V=V^0$
\[
\begin{aligned}
\sin(L+c-i)&=\frac{q}{m}\cos(L+Q)\\
\delta&=m\cos(L+c-i)+q\sin(L+Q)\\
&=\frac{m\cos(Q+i-c)}{\cos(L+Q)}=\frac{q\cos(Q+i-c)}{\sin(L+c-i)}
\end{aligned}
\]

II.\quad For $V=V^0+180^0$
\[
\begin{aligned}
\sin(L+c+i)&=-\frac{q}{m}\cos(L+Q)\\
\delta&=-m\cos(L+c+i)+q\sin(L+Q)\\
&=-\frac{m\cos(Q-c-i)}{\cos(L+Q)}=\frac{q\cos(Q-c-i)}{\sin(L+c+i)}
\end{aligned}
\]

III.\quad The same for $V=V^0+90^0$ and $V=V^0+270^0$
\[
\begin{aligned}
\sin i\cos(L+c)&=-\frac{q}{m}\cos(L+Q)\\
\delta&=m\sin i\sin(L+c)+q\sin(L+Q)\\
&=-\frac{m\sin i\sin(Q-c)}{\cos(L+Q)}=\frac{q\sin(Q-c)}{\cos(L+c)}
\end{aligned}
\]

In our example the values of $L$ in the last two cases differ instead of agreeing, by $3'\ 3''$ and $2'\ 8''$ respectively. These differences are due partly to random errors of observation and partly to several circumstances acting together. The initial magnetic orientation is slightly uncertain. The magnetic declination, and therefore the value of $V^0$, varies over the course of the observations. The horizontal circle has a slight eccentricity that cannot be checked in the absence of a double reading. Finally, setting the needle down with the cups can make the pivot axis perpendicular to the plane of the circle only imperfectly. We reduce the effects of all these circumstances as far as possible by taking the means of the two settings and using the following values


\SourceLeaf{470}\begin{center}
\begin{tabular}{llr}
for $B$ north pole&$.\quad.\quad.\quad.\quad.$&$L=89^0\ 51'\ 49''$\\
for $A$ north pole&$.\quad.\quad.\quad.\quad.$&$L=90\ 10\ 48$
\end{tabular}
\end{center}
Even so, because of these circumstances the result for the setting perpendicular to the magnetic meridian must always be regarded as slightly less reliable than the results for the positions in the meridian itself, where the effect of those causes may be regarded as imperceptible.

\Article{15}
Let the six remaining numbers of the original $32$ be denoted from now on as follows.
\begin{center}
\begin{tabular}{l|l}
Values of $L$&for $V-V^0=$\\\hline
$f,\ f'$&$0$\\
$180^0-g,\ 180^0-g'$&$180^0$\\
$h,\ h'$&$90^0$ and $270^0$
\end{tabular}
\end{center}
The unprimed symbols refer to $B$ as north pole, and the primed symbols to $A$ as north pole. It is clear that, for the positions in the magnetic meridian, $f,f',g,g'$ are all inclinations of the straight line drawn from the south pole tip of the needle to the north pole tip below the northern horizontal line. The first two are for the position in which the marked face of the needle faces east, and the other two for the opposite position. For the positions in the plane perpendicular to the magnetic meridian, $h,h'$ are instead the inclinations of that same straight line to the eastern or western horizontal line, according as the marked face of the needle faces south or north.

Of the elements on which these six quantities depend, $q$ is entirely constant, and $i$ must be taken to be the same for all, since we cannot account for the small fluctuations that may occur over the course of the observations. But $Q,m,c$ change their values after restroking. In particular $Q$ changes by exactly $180^0$, while $m$ and $c$ change without any definite relation to their former values beyond the following. If the poles are reversed by a uniform stroking procedure using str\SourceLeaf{471}ong stroking bars, we can be sure that the difference, and in the case of $c$ also the absolute values, cannot be very considerable. If I now let the unprimed symbols $Q,m,c$ denote the particular values for the observations with $B$ as north pole, and, for the observation with $A$ as north pole, put $Q+180^0,m',c'$ in their place, the general equations of the preceding article become the following six
\[
\begin{array}{lll}
\sin(f+c-i)=\frac{q}{m}\cos(f+Q)&.\ .\ .\ .\ .\ .\ .&(1)\\[4pt]
\sin(g-c-i)=\frac{q}{m}\cos(g-Q)&.\ .\ .\ .\ .\ .\ .&(2)\\[4pt]
\sin i\cos(h+c)=-\frac{q}{m}\cos(h+Q)&\phantom{.}\ .\ .\ .\ .\ .\ .&(3)\\[4pt]
\sin(f'+c'-i)=-\frac{q}{m'}\cos(f'+Q)&\phantom{.}\ .\ .\ .\ .\ .\ .&(4)\\[4pt]
\sin(g'-c'-i)=-\frac{q}{m'}\cos(g'-Q)&\phantom{.}\ .\ .\ .\ .\ .\ .&(5)\\[4pt]
\sin i\cos(h'+c')=\frac{q}{m'}\cos(h'+Q)&.\ .\ .\ .\ .\ .\ .\ .&(6)
\end{array}
\]

\Article{16}
In theory these six equations suffice to determine the six unknown quantities $c,c',\frac{q}{m},\frac{q}{m'},Q,i$. The solution may be given here, although it has no practical value at all, since the enormous effect of the unavoidable observational errors on the final results makes the method quite unusable.

Multiplying equations $1,2,3$ respectively by
\[
\sin(g+h),\quad\sin(f-h),\quad\sin(f+g)
\]
and adding, we obtain after a few easy reductions
\[
\sin(f+c).\sin(g+h)=\sin(g-c).\sin(h-f)
\]
from which $c$ is easily determined, most conveniently by the formula
\[
\operatorname{tang}(c+\tfrac{1}{2}f-\tfrac{1}{2}g)=-\operatorname{tang}\tfrac{1}{2}(f+g)^2.\operatorname{cotang}(h-\tfrac{1}{2}f+\tfrac{1}{2}g)
\]
Similarly, from equations $4,5,6$ we obtain
\[
\operatorname{tang}(c'+\tfrac{1}{2}f'-\tfrac{1}{2}g')=-\operatorname{tang}\tfrac{1}{2}(f'+g')^2.\operatorname{cotang}(h'-\tfrac{1}{2}f'+\tfrac{1}{2}g')
\]
\SourceLeaf{472}The numbers in our example are
\[
\begin{array}{ll}
f=67^0\ 26'\ 11''&f'=67^0\ 58'\ 11''\\
g=67\ 43\ 46&g'=67\ 35\ 35\\
h=89\ 51\ 49&h'=90\ 10\ 48
\end{array}
\]
from which the formulas above give
\[
c=+12'\ 21''\qquad c'=-14'\ 18''
\]
These values are already so large as to be scarcely plausible. Their unreliability is evident when one works out how small errors in the numbers on which they are based affect them. The differential formula for this purpose can be given in several forms. One is the following.
\[
\begin{aligned}
\mathrm{d}c={}&-\frac{\sin(g-c).\sin(h+c)}{\sin(h-f).\sin(f+g)}\mathrm{d}f\\
&+\frac{\sin(f+c).\sin(h+c)}{\sin(g+h).\sin(f+g)}.\mathrm{d}g
+\frac{\sin(f+c).\sin(g-c)}{\sin(h-f).\sin(h+g)}.\mathrm{d}h
\end{aligned}
\]
The same formula holds for $\mathrm{d}c'$ if we simply replace $f,g,h$ by $f',g',h'$. Applied to our calculation, the formulas give
\[
\begin{aligned}
\mathrm{d}c&=-3{,}435\mathrm{d}f+3{,}441\mathrm{d}g+5{,}876\mathrm{d}h\\
\mathrm{d}c'&=-3{,}499\mathrm{d}f'+3{,}494\mathrm{d}g'+5{,}993\mathrm{d}h'
\end{aligned}
\]
If we consider, then, that the values of $h$ and $h'$ are themselves less reliable and may well contain errors of one or a couple of minutes, it is clear that the values found for $c$ and $c'$ deserve no confidence

For completeness, I give next the way in which the other unknown quantities can be found.

Combining equations $(1)$ and $(2)$ gives
\[
\cos i=-\frac{q}{m}.\frac{\sin(f+g)\sin(Q-c)}{\sin(2c+f-g)}\quad.\ .\ .\ .\ .\ .\ .\quad(7)
\]
and hence, using equation $(3)$ as well,
\[
\operatorname{tang}i=\frac{\sin(2c+f-g)}{\sin(f+g).\cos(h+c)}.\frac{\cos(Q+h)}{\sin(Q-c)}
\]
In much the same way equations $4$---$6$ give
\expandafter\def\csname HistoricalNoteText473\endcsname{All the calculations in this paper were carried out with the greatest precision, but the fractions of seconds were left out in printing. Anyone who continues the calculation with the shortened intermediate numbers will therefore sometimes find slightly different results.}
\SourceLeaf{473}\[
\operatorname{tang}i=\frac{\sin(2c'+f'-g')}{\sin(f'+g').\cos(h'+c')}.\frac{\cos(Q+h')}{\sin(Q-c')}
\]
Consequently, if for brevity we write
\[
\frac{\sin(2c'+f'-g').\sin(f+g).\cos(h+c)}{\sin(2c+f-g).\sin(f'+g').\cos(h'+c')}=k
\]
we have
\[
\cos(Q+h).\sin(Q-c')=k\cos(Q+h').\sin(Q-c)
\]
If we put
\[
\begin{aligned}
\cos(h-c')-k\cos(h'-c)&=A\sin B\\
\sin(h-c')-k\sin(h'-c)&=A\cos B\\
\frac{\sin(h+c')-k\sin(h'+c)}{A}&=C
\end{aligned}
\]
this equation takes the simple form
\[
\cos(2Q-B)=C
\]
which determines $Q$. Then $i$ is found from either of the two equations for $\operatorname{tang}i$. Finally, $\frac{q}{m}$ and $\frac{q}{m'}$ are found from $(1)$ or $(2)$ and from $(4)$ or $(5)$. The following points about these calculations are worth noting.

I.\quad To carry out the numerical calculation with precision using the formulas above, $c$ and $c'$ must be computed to much greater accuracy than their absolute unreliability warrants in itself. Otherwise the two determinations of $i,\frac{q}{m},\frac{q}{m'}$ would agree poorly\HistoricalNote{473}. Other, somewhat less simple formulas can be substituted for those above and avoid this drawback. I give them here, omitting the derivation, which is not difficult.
\[
\begin{aligned}
\operatorname{tang}i&=-\frac{2\sin(f+c).\sin(g-c)}{\sin(f+g).\sin(h+c)}.\frac{\cos(Q+h)}{\sin(Q-c)}\\
&=-\frac{2\sin(f'+c').\sin(g'-c')}{\sin(f'+g').\sin(h'+c')}.\frac{\cos(Q+h')}{\sin(Q-c')}
\end{aligned}
\]
\[
k=\frac{\sin(f+g).\sin(f'+c').\sin(g'-c').\sin(h+c)}{\sin(f'+g').\sin(f+c).\sin(g-c).\sin(h'+c')}
\]

\SourceLeaf{474}II.\quad The equation $\cos(2Q-B)=C$, except in the special case $C=\pm1$, always has four distinct solutions, or values between $0$ and $360^0$ for $Q$, differing in pairs by $180^0$. The two values of $Q$ in any such pair correspond to the same value of $i$, but to values of $\frac{q}{m},\frac{q}{m'}$ that agree in magnitude and are opposite in sign. Since these latter quantities must be positive by their nature, one value of $Q$ in each pair is automatically excluded. But if a value of $Q$ gives $\frac{q}{m},\frac{q}{m'}$ opposite signs to each other, the whole pair must plainly be rejected. If that should happen for both pairs, we can conclude only that the observational errors make the combination of equations $1$---$6$ quite unsuitable for determining the unknown quantities. In our example the calculation gives the following. two systems of values
\Needspace{13\baselineskip}
\begin{center}First system\end{center}
\[
\begin{aligned}
Q&=\left\{\begin{array}{rrr}12^0&44'&41''\\192&44&41\end{array}\right.\\
i&=\phantom{\{}\begin{array}{rrr}67&41&33\end{array}\\
\frac{q}{m}&=\mp0.0051395\\
\frac{q}{m'}&=\mp0.0042073
\end{aligned}
\]
\begin{center}Second system\end{center}
\[
\begin{aligned}
Q&=\left\{\begin{array}{rrr}179^0&57'&42''\\359&57&42\end{array}\right.\\
i&=\phantom{\{}\begin{array}{rrr}60&2&11\end{array}\\
\frac{q}{m}&=\mp0.3443905\\
\frac{q}{m'}&=\pm0.3563855
\end{aligned}
\]

Here the second system must plainly be rejected in its entirety, as must the upper value of $Q$ in the first. The value $Q=192^0\ 44'\ 41''$ is thus the only admissible one. It is only by a chance compensation of the observational errors that this gives quite a good value of $i$, and that the ratio of the values of $\frac{q}{m}$ and $\frac{q}{m'}$ deviates so greatly, yet no more greatly, from the ratio of the squares of the swing durations (Article\ 9), to which those values ought to be proportional. Indeed, merely increasing the value of $h'$ by one minute, with the values
\SourceLeaf{475}of the other five quantities $f,g,h,f',g'$ unchanged, produces quite unusable results. The calculation by the method above yields two systems of solutions, in which the inclination is $68^0\ 17'\ 40''$ and $66^0\ 23'\ 12''$ respectively, while in both systems the values of $\frac{q}{m},\frac{q}{m'}$ have opposite signs. This is a striking proof that the calculation must not be based on such combinations.

\Article{17}
If we now set aside the observations in the plane perpendicular to the magnetic meridian, they must either be replaced by other data, or we must make certain arbitrary assumptions that are not strictly correct and content ourselves with the accuracy this affords. In all my observations, a new datum can be obtained from the swing durations determined before and after the poles are reversed. Their squares may be regarded as proportional to the quantities $\frac{q}{m},\frac{q}{m'}$. The same apparatus used to observe these swing durations can also serve to determine $c$ and $c'$ directly. One places the needle in the yoke twice, once with the marked face up and once with it down, observes the positions of the tips relative to the graduated arc, and accounts for any changes in declination by making simultaneous observations with the unifilar magnetometer. The readings that apparatus allows are, however, less precise than \textit{this} application would require (it was not designed for this purpose). But if such an apparatus had much finer graduations, were mounted so that it could not shift, and were read with microscopes, for example, it would indeed be possible to determine $c$ and $c'$ directly with all the precision one could wish for. We would then even have one more datum than necessary, so that a suitable adjustment could further improve the accuracy of the result.

For the time being, then, I replace the missing datum by assuming that the magnetic axis of the needle is unchanged by the reversal of the poles, or that $c'=c$. All observers who have tried to determine the inclination by a more rigorous calculation than the otherwise generally used formula $i=\frac14(f+g+f'+g')$ have made
\SourceLeaf{476}this assumption. There is indeed reason to believe that it is unlikely to be \textit{far} wrong if the stroking is always done with great care, using the same stroking bars and keeping the needle in the same position in a suitably constructed trough. My own experience shows, however, that appreciable differences in the position of the magnetic axis of the needle can occur despite this precaution, and clear signs of them are often recognisable in the reports of other observers as well. (For example, \textsc{Erman}'s observations of 13 October 1829, treated according to his own principles, give a deviation of the magnetic axis of $36'\ 24''$ for one of the needles, whereas at other times it seems to have been very small). Fortunately, under circumstances like those here, even a considerable error in that assumption can have only a very slight effect on the result.

\Article{18}
On this basis the problem can be solved as follows. With equation $(7)$, already used above,
\[
\frac{\cos i.\sin(2c+f-g)}{\sin(f+g)}=-\frac{q}{m}.\sin(Q-c)
\]
I combine the equation obtained in much the same way from $(4)$ and $(5)$, writing $c$ in place of $c'$ and $\frac{\lambda q}{m}$ in place of $\frac{q}{m'}$,
\[
\frac{\cos i.\sin(2c+f'-g')}{\sin(f'+g')}=-\frac{\lambda q}{m}.\sin(Q-c)\quad.\ .\ .\ .\ .\ .\quad(8)
\]
and so
\[
\lambda\sin(f'+g')\sin(2c+f-g)=\sin(f+g).\sin(2c+f'-g')
\]
which determines $c$, most conveniently by means of formula $(9)$
\[
\operatorname{tang}(2c-\tfrac12(g+g'-f-f'))=\frac{\lambda\sin(f+g)-\sin(f'+g')}{\lambda\sin(f+g)+\sin(f'+g')}.\operatorname{tang}(f-g-f'+g')
\]

It follows further from $(1)$ and $(2)$ that
\[
2\cos i.\sin(f+c)\sin(g-c)-\sin i.\sin(f+g)=\frac{q}{m}.\cos(Q-c)
\]
and hence, on combining this with $(7)$,
\[
\operatorname{cotang}(Q-c)=\frac{\sin(f+g)}{\sin(2c+f-g)}.\operatorname{tang}i-\frac{2\sin(f+c).\sin(g-c)}{\sin(2c+f-g)}
\]
\SourceLeaf{477}In much the same way, from $(4)$, $(5)$ and $(8)$ one obtains
\[
\operatorname{cotang}(Q-c)=\frac{\sin(f'+g')}{\sin(2c+f'-g')}.\operatorname{tang}i-\frac{2\sin(f'+c).\sin(g'-c)}{\sin(2c+f'-g')}
\]
If for brevity we write
\[
\begin{array}{ll}
\operatorname{cotang}(f+c)=F&\operatorname{cotang}(f'+c)=F'\\
\operatorname{cotang}(g-c)=G&\operatorname{cotang}(g'-c)=G'
\end{array}
\]
the two equations take the form
\[
\begin{aligned}
\operatorname{cotang}(Q-c)&=\frac{G+F}{G-F}.\operatorname{tang}i-\frac{2}{G-F}\\
\operatorname{cotang}(Q-c)&=\frac{G'+F'}{G'-F'}.\operatorname{tang}i-\frac{2}{G'-F'}
\end{aligned}
\]
from which we finally obtain
\[
\begin{aligned}
\operatorname{tang}i&=\frac{G'-F'-G+F}{G'F-GF'}\\
\operatorname{cotang}(Q-c)&=\frac{G'+F'-G-F}{G'F-GF'}
\end{aligned}
\]
Once $i$ and $Q$ have been found, $\frac{m}{q}$ can be determined from any of equations $1,2,4,5,7,8$.

In our example we have
\[
\lambda=\left(\frac{5{,}87416}{5{,}83555}\right)^2
\]
and the rest of the calculation gives
\[
\begin{array}{rl}
c=&-0^0\ 1'\ 13''\\
i=&67\ 40\ 54\\
Q-c=&145\ 17\ 10\\
Q=&145\ 15\ 57\\
\dfrac{q}{m}=&0.0055111\\[4pt]
\dfrac{q}{m'}=&0.0055843
\end{array}
\]
The values of $h,h'$ \textit{calculated} from these elements are
\[
\begin{array}{rl}
h=&89^0\ 49'\ 30''\\
h'=&90\ 12\ 59
\end{array}
\]
so the observed values differ from these by $+2'\ 19''$ and $-2'\ 11''$.

\SourceLeaf{478}\Article{19}
In the absence of a direct determination of the ratio between $\frac{q}{m},\frac{q}{m'}$, one has to make two arbitrary assumptions in place of one. Observers have used the following two sets of assumptions.

I.\quad Both $c=0$ and $c'=0$ are assumed. This gives us the formula for $i$
\[
\operatorname{tang}i=\frac{\operatorname{cotg}g'-\operatorname{cotg}f'-\operatorname{cotg}g+\operatorname{cotg}f}{\operatorname{cotg}g'\operatorname{cotg}f-\operatorname{cotg}f'\operatorname{cotg}g}
\]
This is the usual procedure when, following \textsc{Mayer}'s example, the needle has deliberately been loaded with a small side weight. This makes the needle settle at quite different points on the graduated circle from those it would reach without the weight. If no substantially different results are obtained, this gives some assurance that no parts of the graduated circle have magnetism of their own. It is advisable, however, to keep the load moderate. With a heavier load, observational errors would have an unduly large effect on the result, and neglecting $c,c'$ would still have an appreciable adverse effect.

II.\quad The assumption here is that $m=m'$ and $c=c'$. This is evidently a special case of the one treated in the preceding article, so its formulas can be applied directly by setting $\lambda=1$. Formula $(9)$ for $c$ then takes a somewhat simpler form
\[
\operatorname{tang}(2c-\tfrac12(g+g'-f-f'))=\frac{\operatorname{tang}\frac12(f+g-f'-g').\operatorname{tang}\frac12(f-g-f'+g')}{\operatorname{tang}\frac12(f+g+f'+g')}
\]
When $c$ is not required and the aim is only to determine $i$, an elegant procedure for the calculation can be found in \textsc{Erman}'s Reise, part 2, volume 2, p.\ 22.

\Article{20}
The relations developed so far between the observations and the inclination and the other elements hold generally, whether the offset of the centre of gravity from the pivot axis is large or small. The latter will always be the case with needles made by a skilled instrument maker, as long as they are not al\SourceLeaf{479}tered by outside causes (for example, rust spots, grinding, removal of the pivots or weights deliberately added). The formulas then admit a very substantial simplification. As long as $\frac{q}{m}$ or $\frac{q}{m'}$ does not exceed $0.03$, the difference between the sines of $f+c-i,\ g-c-i,\ f'+c'-i,\ g'-c'-i$ and the arcs themselves cannot reach one second. Given the moderate degree of precision attainable in observations with the dip circle, one may therefore safely interchange arc and sine even at considerably larger values of $\frac{q}{m},\frac{q}{m'}$. For the four needles of the \textsc{Robinson} dip circle the values lie within much narrower limits. I will therefore use such a shortened procedure for the observations to be reported here, after first applying it to the example we have been considering.

\Article{21}
If for brevity we set
\[
\begin{aligned}
\frac{206265''q\cos Q}{m}&=t\\
\frac{206265''q\sin Q}{m}&=u
\end{aligned}
\]
then, provided that
\[
f+c-i,\quad g-c-i,\quad f'+c'-i,\quad g'-c'-i
\]
are small enough to be interchanged with their sines, equations $1,2,4,5$ of Article 15 take the following form
\[
\begin{aligned}
i&=f+c-t\cos f+u\sin f\\
i&=g-c-t\cos g-u\sin g\\
i&=f'+c'+\lambda t\cos f'-\lambda u\sin f'\\
i&=g'-c'+\lambda t\cos g'+\lambda u\sin g'
\end{aligned}
\]
The five unknown quantities $i,c,c',t,u$ cannot be determined by four equations, but they can be expressed in terms of one quantity that remains undetermined. If we choose $c'-c$ for that purpose, this shows very clearly to what extent we are justified in neglecting it. In each individual case the elimination itself is most conveniently carried out after substituting the numerical values of the observational data.

\SourceLeaf{480}In our example the four equations become
\[
\begin{aligned}
i&=67^0\ 26'\ 11''+c-0{,}3837t+0{,}9234u\\
i&=67\ 43\ 46-c-0{,}3790t-0{,}9254u\\
i&=67\ 58\ 11+c'+0{,}3801t-0{,}9393u\\
i&=67\ 35\ 35-c'+0{,}3862t+0{,}9368u
\end{aligned}
\]
from which one finds by elimination
\[
\begin{aligned}
i&=67^0\ 41'\ 54''-0{,}0006(c'-c)\\
t&=-934+0{,}0002(c'-c)\\
u&=+648+0{,}5369(c'-c)\\
\tfrac12(c'+c)&=-73+0{,}0037(c'-c)
\end{aligned}
\]

This shows that the arbitrary assumption that $c$ and $c'$ are equal makes a reliable determination of $u$ impracticable. It has no appreciable effect on the values of $i$ and $t$, and even in determining the mean of $c$ and $c'$ its effect is only slight.

The mean of the four equations is
\[
i=67^0\ 40'\ 56''+0{,}0009t-0{,}0011u
\]
where the constant term is the simple mean of $f,g,f',g'$, and could therefore have been taken directly as the inclination. This is indeed the usual procedure, and it may always be used without concern in cases where the values of $f,g,f',g'$ show no great differences.

\Article{22}
Before leaving the example we have been considering, I should add that equations $3$ and $6$ admit much the same simplification as the others. We may put
\[
\begin{aligned}
c&=90^0-h+\frac{\cos h}{\sin i}.t-\frac{\sin h}{\sin i}.u\\
c'&=90^0-h'-\frac{\lambda\cos h'}{\sin i}.t+\frac{\lambda\sin h'}{\sin i}.u
\end{aligned}
\]
In the numerical calculation we may safely substitute for $i$ the value
\SourceLeaf{481}$\frac14(f+g+f'+g')$, which puts these equations in the following form in our example
\[
\begin{aligned}
c&=+491''+0{,}0026t-1{,}0810u\\
c'&=-648+0{,}0034t+1{,}0953u
\end{aligned}
\]
The values of $h$ and $h'$ are based on twice as many settings as the values of $f,g,f',g'$. If only the number of settings mattered, each of these equations would therefore have to be given the weight $2\sin i^2$, with the weight of each of the four equations of the preceding article set $=1$. However, for the reasons given above (Article 14), the determinations of $h,h'$ are considerably less reliable. To simplify the calculation, we may therefore take all six equations to have the same weight. Calculating the five unknown quantities from them in this way by the method of least squares gives
\[
\begin{array}{rl}
i=&67^0\ 40'\ 55''\\
t=&-934\\
u=&-211\\
c=&+719\\
c'=&-880
\end{array}
\]
These values satisfy all the equations to within $1''$ and $2''$. This degree of agreement must, of course, be regarded as a matter of chance, since the data themselves are much less reliable. Nor do the values of $u,c,c'$ deserve confidence. In general, at inclinations as large as those in our regions, the data are quite unsuitable for distinguishing those quantities with even moderate reliability.

\Article{23}
After this review of the various methods of calculation, I turn to the main subject and first tabulate the observations made in the manner described in Article 8. I give here only the quantities denoted by $f,g,f',g'$, leaving out the partial results from which they were derived by the procedure given in Articles $11$---$13$. This is partly to save space and partly because the elements to which they are related,
\SourceLeaf{482}owing to frequent changes to the bearings and cups, had different values on different days. Most of the observations were made between 8 and 11 in the morning. On 16, 22 and 25 June and 17 and 20 July, however, they were made between 4 and 6 in the afternoon.

The columns give the designation of the north-pole end of the needle, the values of $f$ and $g$ or of $f'$ and $g'$ according as $B$ or $A$ was the north pole, and the duration of the horizontal swing.

\Needspace{16\baselineskip}
\begin{center}
\textit{Observations with needle} 1.

\begin{tabular}{rr|c|rrr|rrr|l}
May&20&$B$&$67^0$&$11'$&$0''$&$67^0$&$58'$&$46''$&$5''87152$\\
&21&$A$&&$57$&$1$&&$35$&$14$&$5.81508$\\
&22&$A$&&$56$&$29$&&$36$&$45$&$5.82044$\\
&24&$B$&&$16$&$45$&&$45$&$48$&$5.81557$\\
&31&$B$&&$18$&$1$&&$49$&$41$&$5.82075$\\
June&2&$A$&&$53$&$55$&&$33$&$9$&$5.85778$\\
&4&$A$&&$56$&$38$&&$32$&$10$&$5.86442$\\
&5&$B$&&$24$&$13$&&$46$&$44$&$5.83615$\\
July&6&$A$&&$59$&$41$&&$35$&$21$&$5.83716$\\
&7&$A$&&$58$&$7$&&$37$&$51$&$5.83818$\\
&8&$B$&&$20$&$8$&&$44$&$47$&$5.89602$\\
&9&$B$&&$20$&$43$&&$44$&$25$&$5.90035$
\end{tabular}
\end{center}

\Needspace{16\baselineskip}
\begin{center}
\textit{Observations with needle} 2.

\begin{tabular}{rr|c|rrr|rrr|l}
May&20&$A$&$67^0$&$40'$&$57''$&$67^0$&$20'$&$37''$&$5''72416$\\
&21&$A$&&$41$&$8$&&$21$&$5$&$5.72453$\\
&22&$B$&&$43$&$28$&&$50$&$45$&$5.65355$\\
&24&$B$&&$41$&$43$&&$54$&$32$&$5.66875$\\
&31&$A$&&$43$&$34$&&$18$&$29$&$5.67439$\\
June&2&$A$&&$41$&$46$&&$18$&$12$&$5.67665$\\
&4&$B$&&$42$&$42$&&$46$&$57$&$5.68010$\\
&5&$B$&&$44$&$53$&&$50$&$24$&$5.68890$\\
July&17&$B$&&$45$&$20$&&$50$&$17$&$5.70183$\\
&18&$A$&&$40$&$26$&&$22$&$50$&$5.68692$\\
&19&$A$&&$40$&$21$&&$22$&$10$&$5.69677$\\
&20&$B$&&$40$&$40$&&$54$&$19$&$5.66585$
\end{tabular}
\end{center}

\SourceLeaf{483}\Needspace{16\baselineskip}
\begin{center}
\textit{Observations with needle} 3.

\begin{tabular}{rr|c|rrr|rrr|l}
June&8&$B$&$67^0$&$47'$&$58''$&$67^0$&$48'$&$52''$&$6''17149$\\
&9&$B$&&$40$&$55$&&$42$&$28$&$6.18077$\\
&11&$A$&&$30$&$58$&&$32$&$35$&$6.18080$\\
&16&$B$&&$40$&$0$&&$42$&$40$&$6.17046$\\
&18&$B$&&$43$&$13$&&$47$&$40$&$6.18005$\\
&22&$A$&&$27$&$33$&&$39$&$19$&$6.16591$\\
&23&$A$&&$29$&$46$&&$41$&$8$&$6.16948$\\
&25&$A$&&$29$&$3$&&$41$&$7$&$6.17663$\\
July&6&$A$&&$32$&$38$&&$40$&$37$&$6.18305$\\
&7&$B$&&$45$&$56$&&$42$&$12$&$6.17982$\\
&8&$B$&&$46$&$59$&&$43$&$37$&$6.18339$\\
&9&$A$&&$30$&$42$&&$39$&$42$&$6.23905$
\end{tabular}
\end{center}

\Needspace{16\baselineskip}
\begin{center}
\textit{Observations with needle} 4.

\begin{tabular}{rr|c|rrr|rrr|l}
June&8&$A$&$67^0$&$45'$&$9''$&$67^0$&$27'$&$3''$&$5''96200$\\
&9&$B$&&$22$&$56$&$68$&$8$&$28$&$5.91653$\\
&11&$B$&&$23$&$16$&&$7$&$48$&$5.94665$\\
&16&$A$&&$49$&$54$&$67$&$12$&$8$&$6.01785$\\
&18&$B$&&$27$&$48$&$68$&$8$&$45$&$5.93204$\\
&22&$B$&&$26$&$46$&&$3$&$56$&$5.94065$\\
&23&$A$&&$50$&$19$&$67$&$15$&$37$&$5.93939$\\
&25&$A$&&$50$&$4$&&$15$&$22$&$5.94731$\\
July&17&$A$&&$50$&$13$&&$15$&$43$&$5.96850$\\
&18&$A$&&$49$&$57$&&$14$&$48$&$5.96931$\\
&19&$B$&&$22$&$43$&$68$&$9$&$18$&$5.92673$\\
&20&$B$&&$22$&$41$&&$10$&$19$&$5.92783$
\end{tabular}
\end{center}

\Article{24}
In calculating from these observations I shall introduce slightly modified auxiliary quantities in place of $t,u$, used above (Articles 21 and 22). If for one of the needles we denote the duration of a horizontal swing by $n$, the sum of the moments of inertia of the needle and the yoke about the axis of rotation, which is vertical during these swings, by $k$, and the length of the simple seconds pendulum by $l$. then, as is well known,
\[
lmnn\cos i=k
\]
\SourceLeaf{484}Choose a reference swing duration $N$ and a reference inclination about midway among the observed values of $n$ and $i$, respectively, and denote the corresponding value of $m$ by $M$, so that
\[
lMNN\cos I=k
\]
Finally, let
\[
\begin{aligned}
x&=\frac{q\cos Q.\cos I.206265''}{M}\\
y&=\frac{q\sin Q.\sin I.206265''}{M}
\end{aligned}
\]
These quantities are therefore constant for all observations with this needle. The equations then become
\[
\begin{aligned}
i&=f+c-\frac{nn\cos f}{NN\cos I}.\frac{\cos i}{\cos I}.x+\frac{nn\sin f}{NN\sin I}.\frac{\cos i}{\cos I}.y\\
i&=g-c-\frac{nn\cos g}{NN\cos I}.\frac{\cos i}{\cos I}.x-\frac{nn\sin g}{NN\sin I}.\frac{\cos i}{\cos I}.y
\end{aligned}
\]
when $B$ is the north pole. When $A$ is the north pole, we need only reverse the signs of the terms containing $x$ and $y$.

This form has the advantage that the coefficients of $x$ and $y$ always differ little from unity. Indeed, with so small an offset of the centre of gravity in the four needles in question, and with such moderate fluctuations in $n$, we can properly put unity in place of those coefficients. I call this the abbreviated calculation. I nevertheless took the trouble to calculate the $192$ coefficients more precisely, leaving out only the factor $\frac{\cos i}{\cos I}$, although the main benefit of doing so is to show more clearly that the abbreviated calculation is permissible. From now on the unprimed letters $N,x,y$ will refer to needle $1$. The values for the other three needles will be distinguished in order by one, two and three primes. The values chosen for the present calculation are
\[
\begin{aligned}
I&=67^0\ 40'\ 0''\\
N&=5''847785\\
N'&=5.686867\\
N''&=6.181742\\
N'''&=5.949567
\end{aligned}
\]
\expandafter\def\csname HistoricalNoteText485\endcsname{It hardly needs saying that such an error, if it is real at all, can be attributed only to a departure of the pivots from a cylindrical shape. It is constant only insofar as the pivots always rest on the same spots, so at a quite different inclination it could have quite a different value.}
\SourceLeaf{485}To save space, I shall give only as much of the calculations themselves as is needed to follow the procedure in outline. The coefficients that differ most from unity are $0.96895$ and $1.04324$, which occur on 9 and 16 June for needle 4.

\Article{25}
The observations with one needle on each day give two equations. Taking half their sum and half their difference gives two further equations, which we may label I and II. This gives $48$ equations I and just as many equations II. I give the first of each as an example. The original equations from the observations of 20 May with needle 1 are
\[
\begin{aligned}
i&=67^0\ 11'\ 0''+c-1{,}02880x+1{,}00460y\\
i&=67\ 58\ 46-c-0{,}99473x-1{,}01038y
\end{aligned}
\]
and from these we obtain
\[
\begin{aligned}
i&=67^0\ 34'\ 53''-1{,}01176x-0{,}00289y\quad.\ .\ .\ .\ .\ .\quad\text{(I)}\\
c&=+1433''+0{,}01703x-1{,}00749y\quad.\ .\ .\ .\ .\ .\quad\text{(II)}
\end{aligned}
\]
To carry out the test indicated in Article 8, however, I have added a term to equations I by writing $i+e$ in place of $i$, so that $e$ denotes any constant\HistoricalNote{485} error of needle 1. For needles 2, 3 and 4 the presumed constant errors are to be denoted by $e',e'',e'''$.

In this way the $48$ equations I contain altogether $36$ unknown quantities, namely the inclinations on the 24 observation days and the $12$ quantities $x,y,e,x',y',e',x''$, and so on. It must first be noted, however, that the terms containing $y,y',y'',y'''$ all have very small coefficients and are absent altogether from the abbreviated calculation. The largest of these $48$ coefficients is the very $0.00289$ in the example above. If we do wish
\expandafter\def\csname HistoricalNoteText486\endcsname{The observation of 23 September is the one used above as an example in Articles $9$–$22$. The other two are given below in Article $30$.}
\SourceLeaf{486}to allow for the small effect, amounting to only a few seconds, we must first derive the values of $y,y',y'',y'''$ from elsewhere. Even rough approximations are in any case sufficient for this purpose.

\Article{26}
For this derivation only equations II are available to us. But in the $12$ equations of this group that refer to any one needle, the letter $c$ represents different values, since its value can change each time the needle is restroked. It is easy to see, then, that these values of $c$ cannot be eliminated from the equations and that we are therefore \textit{forced} to resort to a somewhat precarious hypothesis. My hypothesis is as follows. Despite all the considerable fluctuations in $c$, using the same stroking procedure throughout will yield a mean value of $c$. I therefore assume that this mean is the same for both arrangements of the poles. Of course, only very imperfect cancellation can be expected if the needle has been restroked only a few times, so the value of $y$ derived in this way will be rather uncertain. This uncertainty is unavoidable unless the values of $c$ are determined with a special apparatus (see Article 17 above). To apply this principle, for each needle we first separate the equations II referring to $B$ as north pole from those for which $A$ was the north pole. We then divide each set into as many groups as there were distinct magnetic states in that set. We take the mean of the equations belonging to each group (if it contains more than one) and then the mean of these partial means. Equating the resulting mean values gives the equation that determines $y$. To illustrate, I give the \textit{abbreviated} calculation for needle 1. \textit{For this purpose} I used, in addition to the $12$ observations above, three further observations\HistoricalNote{486}, on 1 August, 7 August and 23 September. Over the entire peri\SourceLeaf{487}od the needle was restroked nine times, giving ten distinct states, five for each arrangement of the poles.

\Needspace{16\baselineskip}
\begin{center}\textit{Needle} 1, $B$ \textit{north}\end{center}
\[
\begin{array}{rr|l}
&&\qquad c+y=\\\hline
\text{May}&20&.\ .\ .\quad+1433''\\
&\begin{array}{r}24\\31\end{array}&\left.\begin{array}{r}+871''\\+950\end{array}\right\}+910\\
\text{June}&5&.\ .\ .\quad+675\\
\begin{array}{r}\text{July}\\\phantom{\text{July}}\\\text{Aug.}\end{array}&\begin{array}{r}8\\9\\1\end{array}&\left.\begin{array}{r}+739\\+711\\+720\end{array}\right\}+723\\
\begin{array}{r}\phantom{\text{Aug.}}\\\text{Sept.}\end{array}&\begin{array}{r}7\\23\end{array}&\left.\begin{array}{r}+584\\+528\end{array}\right\}+556\\\cline{3-3}
\multicolumn{2}{r}{\text{Mean}}&c+y=+859''
\end{array}
\]
\Needspace{16\baselineskip}
\begin{center}\textit{Needle} 1, $A$ \textit{north}\end{center}
\[
\begin{array}{rr|l}
&&\\\hline
\text{May}&\begin{array}{r}21\\22\end{array}&\left.\begin{array}{r}-653''\\-592\end{array}\right\}-623''\\
\text{June}&\begin{array}{r}2\\4\end{array}&\left.\begin{array}{r}-623\\-734\end{array}\right\}-678\\
\text{July}&\begin{array}{r}6\\7\end{array}&\left.\begin{array}{r}-730\\-608\end{array}\right\}-669\\
\text{Aug.}&\begin{array}{r}1\\7\end{array}&\left.\begin{array}{r}-720\\-785\end{array}\right\}-752\\
\text{Sept.}&23&.\ .\ .\quad-680\\\cline{3-3}
\multicolumn{2}{r}{\text{Mean}}&c-y=-680
\end{array}
\]
from which $y=+769''$ follows. The full calculation gave
\[
\begin{aligned}
\text{with }B\text{ north},\quad c&=+859+0.00102x-1.00290y\\
\text{with }A\text{ north},\quad c&=-680+0.00082x+0.99915y
\end{aligned}
\]
from which
\[
y=+769''+0.00097x
\]
follows. In the same way we find for the other three needles
\SourceLeaf{488}\[
\begin{aligned}
y'&=+456''-0.00192x'\\
y''&=-101+0.00134x''\\
y'''&=+1107+0.00224x'''
\end{aligned}
\]
For needle 1, the fluctuations in the values of $c$ reach $14\frac12$ minutes. For needles 2 and 3 they reach $4\frac12$ minutes, and for needle 4, $10$ minutes. So that no special significance is attached to the fact that the value for needle 1 that differs most occurs on the very first day of observation, I should add that, for this needle and for the others, the poles had already been reversed often before the observations reported here, always with the same care and using the same means of stroking.

\Article{27}
After substituting the values of $y,y',y'',y'''$ in equations I, we are left with $32$ unknown quantities. If we then always subtract the two equations for observations on the same day from one another, we obtain $24$ new equations containing only the eight unknown quantities $x,x',x'',x''',e,e',e'',e'''$. The last four occur only in pairwise differences, so if we set
\[
\begin{aligned}
e'-e&=d'\\
e''-e&=d''\\
e'''-e&=d'''
\end{aligned}
\]
only seven unknown quantities remain. The coefficients of $d',d'',d'''$ in these equations are all $+1$ or $-1$, and the coefficients of $x,x',x'',x'''$ all differ very little from $+1$ or $-1$. To find the seven unknown quantities by the method of least squares, we can safely omit multiplication by the respective coefficients when forming the normal equations for $x,x',x'',x'''$. Simple addition is therefore all that is needed to form all seven normal equations. This gives the following normal equations
\SourceLeaf{489}\[
\begin{aligned}
0&=+4804''+12.00266x-0.00708x'+0.01900x''\\
0&=-5806+0.01559x+12.01005x'-0.00072x'''\\
0&=-3228+0.00145x+12.00544x''+0.04561x'''\\
0&=-5267+0.01786x'-0.00489x''+12.00343x'''\\
0&=-297+0.02717x+0.11088x'-0.04723x'''\\
 &\qquad-12d'+4d'''\\
0&=-241+0.06326x+0.05839x''-0.08085x'''\\
 &\qquad-12d''+8d'''\\
0&=+254-0.02682x'-0.02676x''+0.12808x'''\\
 &\qquad+4d'+8d''-12d'''
\end{aligned}
\]
and from them the values
\[
\begin{aligned}
x&=-400''\\
x'&=+484\\
x''&=+267\\
x'''&=+438\\
d'&=-22\\
d''&=-23\\
d'''&=+1
\end{aligned}
\]
In place of the last three we may also set
\[
\tfrac14(e+e'+e''+e''')=\varepsilon
\]
and write
\[
\begin{aligned}
e&=+11''+\varepsilon\\
e'&=-11+\varepsilon\\
e''&=-12+\varepsilon\\
e'''&=+12+\varepsilon
\end{aligned}
\]
The common part $\varepsilon$ clearly cannot be determined from the available data. Substituting the values found for $x,e,x',e'$ and so on in equations I (from which $y,y'$ and so on have already been removed) and leaving out $\varepsilon$ now gives us the following $48$ inclinations
\SourceLeaf{490}\Needspace{29\baselineskip}
\begin{center}
\begin{tabular}{rr|c|rrr||c|rrr}
&&Needle&&&&Needle&&&\\\hline
May&20&$1$&$67^0$&$41'$&$25''$&$2$&$67^0$&$39'$&$12''$\\
&21&&&$39$&$21$&&&$39$&$31$\\
&22&&&$39$&$51$&&&$39$&$22$\\
&24&&&$37$&$43$&&&$40$&$21$\\
&31&&&$40$&$17$&&&$39$&$17$\\
June&2&&&$36$&$39$&&&$38$&$16$\\
&4&&&$37$&$31$&&&$37$&$0$\\
&5&&&$41$&$56$&&&$39$&$48$\\
&8&$3$&&$44$&$12$&$4$&&$43$&$14$\\
&9&&&$37$&$27$&&&$38$&$15$\\
&11&&&$36$&$27$&&&$38$&$1$\\
&16&&&$37$&$6$&&&$38$&$17$\\
&18&&&$41$&$12$&&&$40$&$48$\\
&22&&&$38$&$5$&&&$37$&$51$\\
&23&&&$40$&$6$&&&$40$&$2$\\
&25&&&$39$&$45$&&&$39$&$49$\\
July&6&$1$&&$40$&$42$&$3$&&$41$&$17$\\
&7&&&$41$&$11$&&&$39$&$49$\\
&8&&&$39$&$5$&&&$41$&$3$\\
&9&&&$39$&$12$&&&$39$&$57$\\
&17&$2$&&$39$&$55$&$4$&&$40$&$7$\\
&18&&&$39$&$56$&&&$39$&$31$\\
&19&&&$39$&$35$&&&$38$&$32$\\
&20&&&$39$&$43$&&&$39$&$1$
\end{tabular}
\end{center}

\Article{28}
The discrepancies between the two determinations of the inclination on each day must now give us a measure of the uncertainty of the observations themselves. The largest discrepancy (on 24 May) is $2'\ 38''$, and the sum of the squares of all $24$ differences, taking the second as the unit, is $124389$. It is easy to deduce from the principles of probability that, if we assign equal reliability to observations with each of the four needles (and there is no reason to depart from this assumption), the mean uncertainty of a value of $\tfrac12(f+g)$ or $\tfrac12(f'+g')$ obtained from the observations and used as the basis of our calculations, as far as can be judged from our figures, must be taken to be
\[
=\BareRoot{\frac{124389}{34}}=60''5
\]
\SourceLeaf{491}insofar as we are dealing only with chance or irregular errors of observation. The mean of two such figures based on independent observations will therefore have a mean unreliability of
\[
=\BareRoot{\frac{124389}{68}}=42''8
\]
and this can also be taken as the mean error of an inclination determined in the usual way (that is, with one needle but in \textit{both} arrangements of the poles), provided that the small correction to $\tfrac14(f+g+f'+g')$ is either considered quite imperceptible or can be based on an already established determination of $u$ or $y$ (see Article 21). It goes without saying that this error estimate applies in the first instance only to this instrument and to observations made under circumstances closely similar to those of the observations on which it rests. With fewer than eight settings in each combination, reliability would be lower, although I would hesitate to claim that the mean error of the final result varies exactly inversely as the square root of the number of settings multiplied by the use of cups. On the other hand, I should point out that throughout the observations above the bearings could not be adjusted as perfectly as I wished and as I actually achieved later by using the apparatus mentioned above (Article 5). The figure above thus already includes any increase in observational errors that may result from imperfect bearing adjustment. An effect of constant magnitude is all the less likely here, since the bearings were altered very often. I therefore have reason to expect that future observations with the same instrument will, if anything, show still smaller errors.

A separate investigation, whose details I pass over here, also showed that the mean uncertainty of the $48$ inclinations given in the preceding article differs little from the mean uncertainty of $\tfrac12(f+g)$, and that the means of each corresponding pair, to be given together in Article 30, should be assigned nearly twice the weight, and hence the mean error $42''8$.

\SourceLeaf{492}\Article{29}
A particularly striking and welcome result is how small the values found for $e,e',e'',e'''$ are, or, more precisely at this stage, how small their differences from their mean $\varepsilon$ are. A separate investigation showed that the weight of these determinations is $\frac{96}{11}$ times that of $\tfrac12(f+g)$, so that their mean uncertainty is $=60''5\BareRoot{\frac{11}{96}}=20''5$. This shows that even the existence of differences between $e,e',e'',e'''$ remains quite doubtful. It is extremely unlikely that four needles should have constant errors of almost exactly the same magnitude. We are therefore justified in assuming that they have either no constant errors or only quite imperceptible ones. It might thus seem almost unnecessary to make further tests by taking advantage of the fact that the axes of two of them can be turned.

Some earlier observations do in fact strengthen this conclusion for one of the needles, namely No.\ 4. On four days from 15 to 19 May, observations had been made with needles 3 and 4 in combinations like those used later from 8 to 25 June, with the sole difference that each partial result was based on four settings instead of eight. The pivots of needle 3 were in the same position as later, but those of needle 4 were in a different position, since they were turned by about a quadrant after 19 May. The observations, written as in Article 23, are as follows

\Needspace{8\baselineskip}
\begin{center}\textit{Observations with needle} 3.

\begin{tabular}{rr|c|rrr|rrr|l}
May&15&$B$&$67^0$&$41'$&$26''$&$67^0$&$44'$&$53''$&$6''16166$\\
&17&$B$&&$43$&$52$&&$45$&$52$&$6.20333$\\
&18&$A$&&$33$&$56$&&$39$&$15$&$6.17781$\\
&19&$A$&&$36$&$8$&&$37$&$8$&$6.19566$
\end{tabular}
\end{center}
\Needspace{8\baselineskip}
\begin{center}\textit{Observations with needle} 4.

\begin{tabular}{rr|c|rrr|rrr|l}
May&15&$A$&$67^0$&$14'$&$28''$&$67^0$&$47'$&$49''$&$5''94332$\\
&17&$B$&$68$&$5$&$39$&&$36$&$36$&$5.92034$\\
&18&$B$&&$3$&$30$&&$36$&$13$&$5.94235$\\
&19&$A$&$67$&$3$&$4$&&$59$&$47$&$5.94663$
\end{tabular}
\end{center}

All the observations were made in the morning.

\SourceLeaf{493}
For needle $3$, the values of $x'',y'',e''$ found above were used in the calculation. For needle $4$, however, the values of $x''',y''',e'''$ had to be derived as well as possible from these observations themselves, giving
\[
\begin{aligned}
y'''&=-1103''\\
x'''&=+556''\\
e'''-\varepsilon&=+24''
\end{aligned}
\]
The determination of $y'''$ is certainly very uncertain, since it is based on so few observations. Its effect on the reduction of $\tfrac12(f+g)$ remains quite insignificant, however, since the largest coefficient of $y'''$ in equations I is only $0.00341$. The results for $i$ are then as follows
\begin{center}
\begin{tabular}{rr||rrr||rrr}
\multicolumn{2}{c||}{}&\multicolumn{3}{c||}{Needle $3$}&\multicolumn{3}{c}{Needle $4$}\\\hline
May&15&$67^0$&$38'$&$57''$&$67^0$&$40'$&$0''$\\
&17&&$40$&$36$&&$41$&$36$\\
&18&&$41$&$15$&&$40$&$15$\\
&19&&$41$&$19$&&$40$&$16$
\end{tabular}
\end{center}
The weight of the determination of $e'''-\varepsilon$ is here only twice that of $\tfrac12(f+g)$. Since the observations themselves are considerably less precise than the later ones, it is clear that the value found here gives as little support for the existence of a constant error as the value derived from the later observations.

The large differences between the values of $x'''$ and $y'''$ and those found above (Articles 26 and 27) prove only that the rotatable part of the needle, considered by itself, does not have its centre of gravity on the pivot axis. This is of little consequence anyway.

\Article{30}
I now give the final results for the inclination from all the observations considered. I also include the results of the observations already mentioned for 1 and 7 August, which were made with needle 1 in exactly the same way as on 23 September. The observations themselves were as follows
\SourceLeaf{494}
\begin{center}
\begin{tabular}{r||rrr||rrr}
&\multicolumn{3}{c||}{August $1$}&\multicolumn{3}{c}{August $7$}\\\hline
$f$&$67^0$&$20'$&$12''$&$67^0$&$22'$&$41''$\\
$g$&&$44$&$11$&&$42$&$8$\\
$f'$&&$59$&$53$&$68$&$1$&$56$\\
$g'$&&$35$&$53$&$67$&$35$&$46$
\end{tabular}
\end{center}
\Needspace{20\baselineskip}
\begin{center}\textit{Inclination determinations}

\begin{tabular}{rrr|rrr||rr|rrr}
1842&May&15&$67^0$&$39'$&$28''$&June&18&$67^0$&$41'$&$0''$\\
&&17&&$41$&$6$&&22&&$37$&$58$\\
&&18&&$40$&$45$&&23&&$40$&$4$\\
&&19&&$40$&$47$&&25&&$39$&$47$\\
&&20&&$40$&$18$&July&6&&$41$&$0$\\
&&21&&$39$&$26$&&7&&$40$&$30$\\
&&22&&$39$&$36$&&8&&$40$&$4$\\
&&24&&$39$&$2$&&9&&$39$&$34$\\
&&31&&$39$&$47$&&17&&$40$&$1$\\
&June&2&&$37$&$27$&&18&&$39$&$44$\\
&&4&&$37$&$15$&&19&&$39$&$4$\\
&&5&&$40$&$52$&&20&&$39$&$22$\\
&&8&&$43$&$43$&Aug.&1&&$39$&$57$\\
&&9&&$37$&$51$&&7&&$40$&$26$\\
&&11&&$37$&$14$&Sept.&23&&$40$&$54$\\
&&16&&$37$&$42$&&&&&
\end{tabular}
\end{center}
The mean of all $31$ determinations, with no distinction between their weights, is
\[
67^0\;39'\;44''
\]
and may be taken as valid for 21 June. The mean of the $24$ determinations from 20 May to 20 July alone, with 19 June as the corresponding mean date, is
\[
67^0\;39'\;31''
\]
\Article{31}
The differences between the inclinations for the individual $31$ days and their mean are made up of the remaining effect of observational errors and the actual variations in the inclination itself. For the
\SourceLeaf{495}individual days these components cannot be separated, but with so many observations in the series it may be worth trying to estimate the mean of the actual fluctuations. With this in mind, I first reduced the inclinations, on the assumption of a regular annual decrease of $3$ minutes, to 21 June, and then added the squares of the differences from the mean. Dividing this sum, $220184$, by $30$ gives $7339.5$ as the square of the mean error to which we would be exposed if we chose one of those $31$ inclinations at random and took it as the mean inclination for the time of the observations. If the unequal reliability of the three groups of observations is to be taken into account, the principles of probability give the following equation, where we denote the mean error for the first four observations by $m'$, that for the last three by $m''$, that for the remaining $24$ by $m$, and the mean fluctuation of the inclination itself by $M$
\[
7339{,}5=\frac{24mm+4m'm'+3m''m''}{31}+MM
\]
For $mm$ we found the value $1829.25$ above, or at least we may regard this number as a sufficient approximation. For the other seven observations, in the absence of a reliable measure, we may take as a basis the number of settings from which the results were derived, and thus set
\[
m'm'=2mm,\qquad m''m''=\tfrac85mm
\]
This gives
\[
MM=7339{,}5-\tfrac{184}{155}.1829{,}25=5168
\]
and $M=71''9$.

\Article{32}
I had already made a series of inclination observations with the same instrument and at the same site in the previous year. I give only the final results here.
\SourceLeaf{496}
\begin{center}
\begin{tabular}{rrrcrrr}
1841&Sept.&22&$.\quad.\quad.$&$67^0$&$40'$&$20''$\\
&&24&&&$40$&$53$\\
&&27&&&$46$&$41$\\
&Oct.&2&&&$42$&$57$\\
&&7&&&$42$&$14$\\
&&10&&&$42$&$40$\\
&&12&&&$43$&$15$\\
&&20&&&$44$&$2$\\
&&20&&&$42$&$5$\\
&&22&&&$42$&$52$\\\hline
Mean,&Oct.&8&&$67^0$&$42'$&$48''$
\end{tabular}
\end{center}
The first eight observations were made in much the same way as this year's. On each day two needles (Nos.\ 1 and 2) were used without reversing the poles between observations. The last two observations were made in the usual way, the second on 20 October with needle 4 and the one on 22 October with needle 3. On 27 September and 10 October the observations were made between $3$ and $5$ in the afternoon, and on all other days in the morning. Each of these $10$ inclinations was based on $16$ settings. For this reason they should be assigned a proportionately smaller weight than the inclinations of $1842$, which were based on $32$, $64$ and $40$ settings respectively.

\Article{33}
All the inclinations given so far still require a small common correction for the effect at the observation site of the bar magnets of the magnetometers in the astronomical observatory and in the magnetic observatory. To remove this effect from the results, $5''15$ must be subtracted throughout (see \textit{Resultate} for 1840. II. Article 6) [p.\ 433 of this volume]

The \textit{absolute} reliability of the inclination determinations still depends on the correctness of the assumption that the instrument itself contains no parts that can exert a magnetic effect on the needle. The instrument I used gives no grounds for such a concern. Some observations I made with a weighted needle in the manner mentioned in Article 18 always showed only discrepan\SourceLeaf{497}cies of a couple of minutes. These can readily be explained by the unavoidable chance errors of observation and the actual anomalies of the inclination itself. The sufficiently close agreement of the values found in Article 11 for the quantity there denoted by $\alpha$ also weighs against the presence of such disturbances. Tests of this kind are of course unsuitable for detecting very small effects, and I must therefore leave further testing to more searching methods.

\Article{34}
To close, I bring together my results and some earlier determinations.
\[
\begin{array}{r|l}
\begin{array}{rrr}1805&\text{Dec.}&\\1826&\text{Sept.}&\end{array}&
\left.\begin{array}{rrr}69^0&29'&\\68&29&26''\end{array}\right\}\ \text{\textsc{von Humboldt}}\\
\begin{array}{rrr}1837&\text{July}&1\\&\text{—}&\text{—}\end{array}&
\left.\begin{array}{rrr}67&47&0\\67&53&30\end{array}\right\}\ \text{\textsc{Forbes}}\\
\begin{array}{rrr}1841&\text{Oct.}&8\\1842&\text{June}&21\end{array}&
\begin{array}{rrr}67&42&43\\67&39&39\end{array}
\end{array}
\]
I took the first two observations from the \textit{Additions} to volume XIII of the \textit{Voyage aux régions équinoxiales} (p.\ 152). The first was made with a dip circle by \textsc{Lenoir}, the second with an instrument by \textsc{Gambey}. The second is based on observations with two needles. Their results are given there as $68^0\;30'\;7''$ and $68^0\;28'\;15''$, with which the mean given in the same place does not agree. A printing error has probably made the figure for the second needle $30''$ too small. I do not know where the observation in $1805$ was made. In $1826$ the site was in an open field a few hundred paces east of the observatory.

\textsc{Forbes}'s observations are printed in the \textit{Transactions of the Royal Society of Edinburgh}, vol.\ XV, part 1, pp.\ 31 and 32. They were made in the observatory garden with a \textsc{Robinson} instrument smaller than the one here, using two needles $6$ English inches long. The observer himself regards the second needle as the better one.

I chose to leave out the observations that \textsc{Mayer} made in March 1814 and that the \textit{Commentationes recent. Soc. Gotting.}, vol.\ III, pp.\ 36 and 37, re\SourceLeaf{498}ports, since they deserve no confidence at all. The test he himself gives on p.\ 35 shows how imperfect the instrument \textsc{Mayer} used was. With the instrument in a fixed position, ten repeated settings gave differences of more than one degree. His results for the inclination itself, obtained on two different days, differ by half a degree.

My own observation of 23 June 1832, given in \textit{Intens. vis magneticae terrestris}, Article 27, equally deserved no place here, both because of the imperfection of the instrument and because of the room in the observatory. The ironwork nearby was bound to affect the result considerably and can be shown to have necessarily increased the inclination.

\Needspace{11\baselineskip}
The inclinations given can be reconciled very well by assuming a uniform annual decrease of $3$ minutes, or more precisely $3'\;2''3$, if we use the result of the second needle in \textsc{Forbes}'s observations. The only discrepancies left may properly be attributed to observational errors and fluctuations in the inclination acting together. According to \textsc{Hansteen}'s investigations of observations at other European sites, however, the annual decrease has gradually slowed. We should therefore take the figure given only as a mean value valid around $1829$, and look to future observations to confirm the departure from uniformity and determine it more precisely.


\par\endgroup
